\subsection{Guarded Scripts: The Original Heavy-Meta Counter}
\label{sec:guards}

We now return from interpreter construction to a useful change in
application dispatch. A structural pattern can bind \code{n}, but it does not establish that
\code{n} is a positive integer. A script's \code{where} clause supplies
an additional acceptance test, called a \emph{guard}. This is part of
script selection, not an \code{if} inside an already selected body.
If the guard is false, the metaobject tries the next script for the
same request. The user's guide describes this facility as a library
extension~\cite{guide}; it does not add guards to the ordinary
\code{match} form.

The original example is inside the closing block comment of
\filepath{lib/heavy-meta.abclr2}, not in a separate example file:

\wholelisting[firstline=133,lastline=144]{lib/heavy-meta.abclr2}

Loading the library defines the generator \code{heavy-meta}, not this
commented counter. In a fresh session, load the library before entering
the counter definition above. Then replay its original commands:

\begin{codeblock}
\begin{lstlisting}
(inc-lib "lib/heavy-meta" :use-source t)
; Enter the original counter definition after loading the library.
[counter <= [:add 123]]
[counter <= [:add -100]]
[counter <== [:value]] ; => 123
\end{lstlisting}
\end{codeblock}

The two add requests are past-type sends. Neither add script replies;
do not replace them by now-type sends expecting an acknowledgement.
The value request is now-type and explicitly replies with \code{c}.
The reset request likewise changes state without replying.
The automated \code{guards-check} extracts original lines 133--144
byte-for-byte and replays lines 146--148 separately, leaving the
selected code unchanged.

\paragraph{Trace both structural success and guard failure.}
For \code{[:add 123]}, the first pattern binds \code{n = 123}.
\code{integerp} is true and \code{plusp} is true, so the body adds 123
to \code{c}. For \code{[:add -100]}, that same pattern still matches
and binds \code{n = -100}. Only the guard fails. The metaobject then
tries \code{[:add \_]}, whose wildcard accepts the argument without
binding it. That script has no body and intentionally ignores the
request, but its execution still finishes the dispatch cycle.

The \code{and} evaluates left to right and stops at the first false
result. Consequently an argument such as \code{"7"} fails \code{integerp}
without asking \code{plusp} to compare a string. Zero is an integer
but not positive, so it also reaches the fallback. Lisp's \code{nil}
is false; other values, including zero, are true when used directly
as guard results. A guard need not return exactly \code{t}.

For \code{[:value]}, neither add pattern matches; their guards are
not evaluated at all. The value script has no guard and replies
immediately once selected. Ordering still matters: placing the
unguarded add fallback first would accept every add request before
the guarded script could be considered.

\paragraph{Ensure that the object actually uses heavy-meta.}
The compiler records the guard separately from the body. In the
ordinary counter definition, it selects \code{heavy-meta} automatically
because neither a metaobject generator nor a group is specified.
The current \filepath{src/translate.lsp} makes that condition explicit.
The guide's shorter description of automatic selection should not be
read as overriding any custom group or metaobject configuration.

If you specify a group, explicitly select a compatible generator:

\wholelisting[firstline=26,lastline=32]{tutorial/examples/09-guards.abclr2}

Loading the library alone does not replace the generator used by every
group. A diagnostic object with an explicit ordinary group but no
\code{meta-gen} accepts its first structural match even when its stored
guard is \code{nil}; the ordinary metaobject does not perform this
guard protocol. Likewise, the non-reifying object translator warns
that a \code{where} field is ignored. Do not use a non-reifying object
as a guarded receiver. If another reflective generator is needed,
it must implement compatible guard-aware selection; naming both
features does not compose their algorithms automatically.

\paragraph{Read the library's continuation chain.}
The heavy metaobject retains the familiar ordinary queue, application
state environment, scripts, evaluator, group manager, and active/dormant
mode. Its \code{:begin} dequeues one request as \code{MRS}, a triplet
of message, reply destination, and sender. Selection is then a series
of asynchronous requests and result continuations:

\begin{enumerate}
  \item Ask the candidate script \code{[:match . MRS]}. Structural
  failure immediately moves on to the remaining scripts.
  \item On success, the result contains \code{[Bindings ScriptBody]}.
  Ask \code{env-gen} for a frame containing those bindings, with the
  object's persistent state as its outer environment.
  \item Ask the script for \code{:where-exp}. A \code{nil} reply means
  there is no guard. Otherwise this reply is a compiled expression
  representation, not its boolean result.
  \item Evaluate that representation with \code{[:do ...]} in the new
  environment. Its result continuation either starts \code{ScriptBody}
  in that same environment or tries the remaining scripts.
  \item Attach \code{[cont \_ [Me <= :end]]} to body execution.
  Completion advances the ordinary queue or restores dormant mode.
\end{enumerate}

The central decision is visible in the original implementation:

\wholelisting[firstline=74,lastline=89]{lib/heavy-meta.abclr2}

Its two \code{if} tests answer different questions: whether a guard
representation exists, and whether evaluating it accepted the request.
An explicitly false guard such as \code{where nil} still has a non-nil
representation; it must be evaluated and rejected, not treated as absent.
The completion continuation is also separate from the script's
application reply. The empty fallback can finish execution without
sending anything to the original caller.

The library first presents a synchronous sketch under
\code{\#+COMMENT}, then the active continuation-passing version.
Continuation-passing style, abbreviated CPS in its comments, makes
each next step a handler of the preceding result. Read the active
chain as the implementation; the sketch is explanatory, not an
additional code path executed before it.

\paragraph{Which bindings can the guard see?}
The guard runs after pattern matching and environment creation. It
can use pattern variables, persistent state, outer lexical bindings,
and the bound reply/sender fields of its script header. The supplemental
example checks the last two explicitly:

\wholelisting[firstline=15,lastline=24]{tutorial/examples/09-guards.abclr2}

In the first driver request, \code{Me} denotes the driver, which is
also the sender received as \code{S}. The guard succeeds and the body
confirms that \code{R} is a non-nil reply destination. The other request
passes \code{nil} as the expected sender, so the guard fails and the
fallback replies \code{nil}.

Body-local \code{temporary} bindings are different: their initialization
belongs to body execution, after guard selection. Do not refer to those
bindings as though they were available before the body starts.
The parser's header order is pattern, optional \code{@ R}, optional
\code{from S}, then optional \code{where expression}.

\paragraph{A failed guard does not undo its effects.}
The supplemental \code{guarded-choice} increments \code{checks} inside
its guard before testing the argument against state binding \code{limit}.
Requests for 5, 11, and the symbol \code{wrong} return
\code{:rejected}, \code{:accepted}, and \code{:rejected}, while
\code{checks} becomes 3. Even failed candidates have executed their
guards. The example demonstrates the mechanism, not a recommendation
to hide side effects in acceptance tests.

\wholelisting[firstline=4,lastline=13]{tutorial/examples/09-guards.abclr2}

Prefer simple predicates. A guard that waits on an ordinary request
to its own receiver can deadlock: that request cannot start while
the current guard is still selecting a script. Other objects can also
change shared state; guard evaluation is not a multi-object transaction.

Reflective script lookup is not the same as guard-aware dispatch.
The heavy metaobject's \code{[:find-script Message]} query calls
\code{search-script}, which checks structural matches. In the original
counter it returns the first add script even for \code{[:add -1]}.
Its \code{:where-exp} remains available for inspection, but the lookup
query does not certify that evaluating that guard would accept the request.

\paragraph{Evidence and exercise.}
Run \code{make -C tutorial guards-check}. Twenty-four checks cover the
original counter, script/guard inspection, state and sender/reply bindings,
failed-guard effects, explicit-group setup and non-reifying restrictions.
Translations, extracted sources and transcripts stay in
\filepath{build/guard-tests/}. The runnable supplemental definitions are in
\filepath{tutorial/examples/09-guards.abclr2}.

Why does \code{[:add -100]} not raise an error or block later value queries?
Its structural match succeeds, its guard fails, and the empty fallback
completes dispatch. Why not put the same test inside the first body's
\code{if}? That would select the first script regardless of the test;
it would not ask dispatch to try the next script. Why must a guard's
side effects be considered separately from acceptance? Evaluation has
already happened when its false result selects another candidate.

\paragraph{Which Heavy-Meta library should you load?}
Use \filepath{lib/heavy-meta.abclr2} for guarded programs. The distribution
also contains \filepath{lib/heavy-meta-old.abclr2}, a historical variant
whose guard protocol does not match the current translator. It is not
needed by the original programs in \filepath{examples/} and is not an
interchangeable alternative. Both files define \code{heavy-meta} and
provide the feature \code{heavy-meta}; do not load both into the same
session. The old file is historical reference material, not another
step in this book's learning path.
